\sectionguard
\section*{The Propers: Notable and Quotable}

Five afterlives of wording from the scriptural propers, each a use that moves
the phrase into a register its apostle or evangelist did not occupy.

\begin{afterlife}{Epistle, Eph. 3:16 --- ``the inner man,'' now the belly}
\textbf{Appointed text and locus.} Marginal no.~1594: \latin{ut det vobis
secúndum divítias glóriæ suæ, virtúte corroborári per Spíritum eius in
interiórem hóminem.}

\textbf{Later texts and loci.} Jerome K. Jerome, \work{Three Men in a Boat (To
Say Nothing of the Dog)} (London: J. W. Arrowsmith, 1889), \textbf{chapter
XIX}, the supper after the Alhambra: ``the odour of Burgundy, and the smell of
French sauces, and the sight of clean napkins and long loaves, knocked as a very
welcome visitor at the door of our inner man.'' Anthony Trollope, \work{Barchester
Towers} (London: Longman, Brown, Green, Longmans, \& Roberts, 1857),
\textbf{chapter VII, ``The Dean and Chapter Take Counsel''}: ``When I say up rose
the archdeacon, I speak of the inner man, which then sprang up to more immediate
action, for the doctor had bodily been standing all along with his back to the
dean's empty fire-grate.''

``Inner man'' is the sole occurrence of that collocation in the Authorized
Version, and it stands at this verse. \textbf{The idiom's ancestry is Pauline
in general rather than this verse's in particular}: the Authorized Version, and
the Douay with it, reads ``the inward man'' at Rom.~7:22 and at 2~Cor.~4:16.
\textbf{At this verse the Douay also has ``the inward man'' and does not carry
the phrase at all}, so the English that carries it here is the Authorized
Version's.

\textbf{The turn.} The words Paul uses for the self strengthened by the Spirit
are now most familiar as a polite word for the belly, and \textbf{that gastric
sense is an idiom of the language and not a joke either novelist invented}: it
stands in 1882, seven years before Jerome K.~Jerome, in Henry William Herbert's
``The Yorkshire Moors,'' \work{Life and Writings of Frank Forester}, vol.~2,
p.~35 --- ``we should be there to refresh the inner man, and take up the fresh
dogs for the afternoon.'' Jerome's narrator has been ten days on cold meat and
jam, and Trollope applies the same manoeuvre to an archdeacon at a chapter
meeting, the Pauline interior self standing up while the body does not.
\end{afterlife}

\begin{afterlife}{Gospel, Luke 14:5 --- ``the ox is in the ditch''}
\textbf{Appointed text and locus.} Marginal no.~1597: \latin{Cuius vestrum
ásinus, aut bos in púteum cadet, et non contínuo éxtrahet illum die sábbati?}

\textbf{Later texts and loci.} \work{San Antonio Daily Light} (San Antonio,
Tex.), \textbf{Thursday 7 October 1897, page 8}, an editorial paragraph headed
``THE EMERGENCY MET'': ``When the ox is in the ditch its no time to stop and
discuss parliamentary ethics; the thing to do is to go to work and pull the ox
out of the the mire''; the 1897 page sets \emph{the the}, and the budget was
stalled while two aldermen figured on parliamentary tactics. Clifton
Johnson, \work{Highways and Byways of the South} (New York: The Macmillan
Company, 1904), chapter ``The Birthplace of Lincoln,'' \textbf{p.~169}, of a
farm wife near Hodgenville, Kentucky: ```The ox is in the ditch,' said she,
referring to the New Testament excuse for Sabbath work where the need is great,
`and Billy must mend this gate if it is Sunday. It was broken yesterday, and I
couldn't sleep last night for thinkin' the hogs might come in and turn over my
kittle of soft soap.'''

\textbf{Johnson's own gloss documents the dependence and the idiom is not the
verse's wording}: the saying has a ditch where the Latin has \latin{púteum} and
both the Douay and the Authorized Version have a pit, so it is a paraphrase of
this verse and never a quotation of it. \textbf{The rival ancestors carry one
half of it each}: Matt.~12:11 has a sheep in a pit on the sabbath but no ox,
and Deut.~22:4, \latin{Si videris asinum fratris tui aut bovem cecidisse in
via}, has the fallen ox but no sabbath. Neither of the nearest Exodus laws
carries that fallen ox: Exod.~23:4 sends back an ox that is straying, and
Exod.~23:5 relieves an ass lying under its burden. Luke~14:5 is the only place
in the canon carrying both, which is why the idiom's sabbath sense points
here.

\textbf{The turn.} The emergency shrinks from a drowning beast to a soap kettle
and then to a stalled city budget, and the verse survives as the standing
American excuse for suspending a rule. It runs through the Library of
Congress's \work{Chronicling America} from 1897 into the 1950s, in Texas,
Georgia, Pennsylvania, New Jersey, Arkansas, Mississippi and Washington, D.C.,
beside a devotional life of its own: the 1942 \work{Evening Star} of Washington
carries it as a Lutheran sermon title in a church-notices column.
\end{afterlife}

\begin{afterlife}{Gospel, Luke 14:10 --- ``Sit Down in the Lowest Room.'' as the title of a secular poem}
\textbf{Appointed text and locus.} Marginal no.~1597: \latin{Sed cum vocátus
fúeris, vade, recúmbe in novíssimo loco: ut, cum vénerit qui te invitávit, dicat
tibi: Amíce, ascénde supérius.}

\textbf{Later text and loci.} Christina G. Rossetti, \textbf{``Sit Down in the
Lowest Room.''}, \work{Macmillan's Magazine}, ed.\ David Masson, \textbf{No.~53,
vol.~IX} (London: Macmillan and Co., \textbf{March 1864}), \textbf{pp.~436--439},
first line ``Like flowers sequestered from the sun.'' Collected as ``The Lowest
Room'' in \work{The Poetical Works of Christina Georgina Rossetti}, ed.\ William
Michael Rossetti (London: Macmillan and Co., 1904), poem at \textbf{pp.~16--19},
editorially dated 30~September 1856, with the editor's note at pp.~460--461.

\textbf{The link is verbatim and it is carried by the title}: the 1864 heading is
the Authorized Version's imperative clause, full stop included, set over a
secular poem, and ``lowest room'' occurs in that Bible at Luke~14:9 and 14:10
only, both inside this appointed pericope. The poem's body reads ``Content to
take the lowest place,'' which is the Douay wording, so the title and the body
draw on two English versions and the appointed Latin lies behind both.

\textbf{The turn.} A rule of banquet precedence and a promise of eschatological
reversal becomes the name for Victorian women's enforced self-effacement, with
the reversal deferred to the trumpet. The poem is a dialogue between two
sisters, one plain and bookish and stung by Homer's heroes, the other beautiful
and married; twenty years on the speaker sits alone --- ``Not to be first: how
hard to learn / That lifelong lesson of the past\dots\ So now in patience I
possess / My soul year after tedious year, / Content to take the lowest place, /
The place assigned me here'' --- and ends ``Yea, sometimes still I lift my heart
/ To the Archangelic trumpet-burst, / When all deep secrets shall be shown, /
And many last be first.''

\textbf{The editor's own note records what the title cost her}: the Gospel
imperative displaced the manuscript's title,
\work{A Fight over the Body of Homer}, ``perhaps the better title of the two'';
Dante Gabriel Rossetti wrote in 1875 of ``a real taint, to some extent, of
modern vicious style\dots\ what might be called a falsetto muscularity'' in
``the long piece called \work{The Lowest Room}''; ``Christina, on receiving this
letter, did not acquiesce in its purport,'' and ``she always retained \work{The
Lowest Room} in succeeding editions.'' Rossetti also wrote a short devotional
lyric ``The Lowest Place'' and a separate devotional piece headed ``Sit down in
the Lowest Room''; \textbf{neither is this poem.}
\end{afterlife}

\begin{afterlife}{Epistle, Eph. 3:18 --- ``the depth and breadth and height'' of a love for one man}
\textbf{Appointed text and locus.} Marginal no.~1594: \latin{ut possítis
comprehéndere cum ómnibus sanctis, quæ sit latitúdo, et longitúdo, et
sublímitas, et profúndum: scire étiam supereminéntem sciéntiæ caritátem
Christi.}

\textbf{Later text and locus.} Elizabeth Barrett Browning, \work{Sonnets from
the Portuguese}, \textbf{Sonnet XLIII}, first published in \work{Poems} (London:
Chapman and Hall, \textbf{1850}), lines 1--4: ``How do I love thee? Let me count
the ways. / I love thee to the depth and breadth and height / My soul can reach,
when feeling out of sight / For the ends of Being and ideal Grace.''

\textbf{The identification is a scholar's and the relation is an echo.} John S.
Phillipson, ``\,`How Do I Love Thee?' --- an Echo of St.~Paul,'' \work{The
Victorian Newsletter}, ed.\ William E. Buckler, \textbf{No.~22 (Fall 1962),
p.~22, note~7}, writes that Barrett Browning ``seems to echo St.~Paul,
Ephesians, III, 17--19, in her famous declaration,'' and closes that ``Sonnet 43
echoes St.~Paul's thought and phraseology while adapting them to a new
context.'' \textbf{Barrett Browning names no source in the sonnet, and the
dependence is Phillipson's identification rather than her own claim.} The sonnet has three of the four
nouns, reordered, and drops \emph{length}, against the Authorized Version's
``the breadth, and length, and depth, and height'' and the Douay's ``the breadth
and length and height and depth'': a reuse of the vocabulary and of the
measuring gesture, not a quotation.

\textbf{The turn.} Paul prays that the Ephesians may take the measure of a love
that passes knowledge; Barrett Browning measures her love for Robert Browning,
keeping the religious lexis --- soul, Grace, faith, saints --- around it, and
from that sonnet the phrase passed into the standard wedding reading and the
greeting-card formula. \textbf{Phillipson reads the transfer as sacralising the
profane rather than secularising the sacred}, the religious lexis moving toward
the beloved and not away from God.
\end{afterlife}

% The gallery's last entry is also its tallest, and on the page it began from
% only its header bar and three lines of the first field fitted, so the box
% opened at the foot of one page and closed near the top of the next. Set at
% the entry's own height, the guard starts it on the following page, where it
% stands whole as every other entry of the gallery does. The section still
% occupies the same three pages.
\Needspace{27\baselineskip}
\begin{afterlife}{Gospel, Luke 14:11 --- ``he that humbleth himself,'' improved}
\textbf{Appointed text and locus.} Marginal no.~1597, the sentence the pericope
ends on: \latin{quia omnis, qui se exáltat, humiliábitur: et qui se humíliat,
exaltábitur.} Douay: ``Because every one that exalteth himself shall be humbled:
and he that humbleth himself shall be exalted.''

\textbf{Later text and locus.} Friedrich Nietzsche, \work{Menschliches,
Allzumenschliches: Ein Buch für freie Geister}, Erster Band, Zweites Hauptstück
(``Zur Geschichte der moralischen Empfindungen''), \textbf{aphorism 87,
entire}: ``Lucas 18,14 verbessert. --- Wer sich selbst erniedrigt, will erhöhet
werden.'' In Helen Zimmern's English of 1910, \work{Human, All-Too-Human: A Book
for Free Spirits}, Part~I, \textbf{printed p.~88}: ``St.~Luke xviii.~14,
Improved. --- He that humbleth himself \emph{wishes to be} exalted.''

\textbf{The English clause is verbatim and its alteration exact.} Zimmern leaves
the Authorized Version's ``He that humbleth himself'' untouched and changes only
\emph{shall be} to \emph{wishes to be}. Nietzsche's German follows the received
Luther wording of Luke~18:14 and turns it at the modal, \textbf{and no
public-domain German Bible stands behind that comparison here}. And the title
claims the act: \emph{verbessert}, improved.

\textbf{Nietzsche's heading cites \emph{Lucas 18,14}, the Pharisee and the
Publican, and not Luke~14:11, the verse this Mass appoints.} The two carry the
identical clause in the Clementine Latin the missal prints, in the Douay and in
the Authorized Version, so the sentence he rewrites is word for word the
sentence the appointed Gospel ends on --- but the parallel is the one he names.
Matt.~23:12 carries the same reversal in different words; the two Lucan verses
are the two that are verbally identical.

\textbf{The turn.} The promise spoken over a dinner table where guests were
scrambling for the best couch is answered eighteen centuries later by a
philosopher who \textbf{agrees that the humble will be exalted and says that is
exactly why they humble themselves} --- the parable's cure re-described as the
disease, and set under a title that claims to be repairing Scripture. It stands
against Bede at the same verse, for whom the humbling and the exalting are the
Lord's acts and not men's, and against Benedict, who makes the sentence the head
of his chapter on humility.
\end{afterlife}

\clearpage
